% UNIFIED INTRODUCTION - Issue #2 Revision
% Date: February 13, 2026
% Follows: Motivation → Solution → Validation narrative
% Emphasizes: Integrated contribution, not separate pieces

\section{Introduction}
\label{sec:intro}

The deployment of Large Language Models (LLMs) presents a fundamental challenge: \textbf{how do we route user queries across models with vastly different costs and capabilities?} Naively routing all traffic to expensive "frontier" models (e.g., GPT-4) wastes resources, while static routing policies fail to adapt to distribution shifts and new model releases.

% ============================================================================
% PART 1: THE PROBLEM (Three Challenges)
% ============================================================================

\subsection{Three Interconnected Challenges}

\paragraph{Challenge 1: Expensive Models Aren't Always Better}
Recent evidence reveals a surprising phenomenon: cheaper models can outperform expensive ones on specific tasks. We discover an \textbf{``Alignment Tax''} where GPT-4-Turbo ($\$10$/1M tokens) actually performs \emph{worse} than Mixtral ($\$0.50$/1M tokens) on $17.6\%$ of prompts, plausibly due to RLHF over-optimization for conversational helpfulness (though the causal mechanism has not been ablated; see Section~\ref{sec:motivation}). Under stylized assumptions (1M prompts/day, current pricing), this creates a routing opportunity on the order of \$2M/year---but only if we can \emph{learn} which prompts favor which models.

\paragraph{Challenge 2: Training Data Doesn't Match Deployment}
Distribution shift between training and deployment data is substantial (PSI=$0.275$, $p<10^{-37}$). When warmup priors trained on historical data are deployed on production traffic with domain mismatch, they can \emph{catastrophically fail}. Static routers like RouteLLM~\cite{ong2024routellm} and FrugalGPT~\cite{chen2023frugalgpt} have no mechanism to detect or recover from such failures.

\paragraph{Challenge 3: New Models Release Monthly}
The LLM landscape is non-stationary: GPT-4o, GPT-5, and competitors release frequently. Standard contextual bandits require thousands of exploratory requests to learn new model capabilities (the ``cold start'' problem), making rapid adoption impractical.

\textbf{The Key Insight:} These challenges are \emph{interconnected}. Solving one without the others leaves production systems vulnerable. For example, learning semantic structure (Challenge 1) doesn't help if priors fail (Challenge 2) or if we can't adopt new models (Challenge 3).

% ============================================================================
% PART 2: OUR SOLUTION (Integrated Approach)
% ============================================================================

\subsection{Our Integrated Approach}

We present \textbf{banditGPT}, a contextual bandit framework that addresses all three challenges simultaneously through an integrated system combining:

\paragraph{1. Semantic Structure Discovery}
We demonstrate that task difficulty in LLM routing relates to \emph{model alignment failures}, not just reasoning complexity. Using principal component analysis on prompt embeddings, we discover that the first principal component (PC1, capturing $3.10\%$ variance) separates tasks by alignment strictness. This semantic structure makes routing \textbf{learnable}—but only if we can safely learn from it.

\paragraph{2. Corralling Meta-Learning} 
To provide safety against prior mismatch, we employ Corralling~\cite{agarwal2017corralling}, a meta-algorithm that adaptively combines two experts: one with warmup priors (semantic knowledge) and one starting fresh (tabula rasa). Through exponential weight updates, Corralling automatically detects when priors fail and switches to cold-start exploration. This provides \textbf{safety guarantees}: $44.3\%$ improvement when priors are harmful, while achieving near-optimal performance ($1.3\times$ vs optimal) when they succeed.

\paragraph{3. Semantic Transfer}
To eliminate cold-start penalties for new models, we introduce \textbf{semantic transfer}: new models inherit preferences ($\theta$) from semantically similar models while resetting confidence ($A$), enabling immediate exploitation of transferred knowledge with maintained exploration flexibility. Statistical validation ($+0.62$ reward improvement on new model release, $p<10^{-7}$) confirms this eliminates cold-start costs.

\paragraph{Why Integration Matters}
\textbf{Alone, each mechanism is insufficient:}
\begin{itemize}
    \item Semantic structure without safety → catastrophic failure on distribution shift
    \item Corralling without transfer → cold-start penalties on new models  
    \item Transfer without meta-learning → no recovery mechanism when transfer fails
\end{itemize}

\textbf{Together, they create a production-ready system} with safety, performance, and adaptability.

% ============================================================================
% PART 3: VALIDATION & CONTRIBUTIONS
% ============================================================================

\subsection{Contributions}

Our work makes three contributions, validated through comprehensive experimentation:

\begin{enumerate}
    \item \textbf{Alignment Tax Discovery:}
    First empirical demonstration that RLHF optimization creates performance \emph{inversions}---$17.6\%$ of prompts where Mixtral outperforms GPT-4 ($p<10^{-143}$, Cohen's $d=1.90$). Structure validated at $594{,}199$ prompts ($317\times$ scale).
    
    \item \textbf{Corralling for LLM Routing:}
    Introduction of Expert Corralling~\cite{agarwal2017corralling} to the routing domain, providing $44.3\%$ improvement vs harmful warmup priors with multi-seed validation (N=$10$ seeds, $75$ ablation configurations).
    
    \item \textbf{Integrated Production System:}
    $68.5\%$ gap closure to oracle (vs $46.2\%$ for RouteLLM) with catastrophic failure detection ($3$--$50$ steps), zero-shot model adoption ($+0.62$ reward, $p<10^{-7}$), and multi-seed statistical testing (N=$10$--$30$) exceeding field standards.
\end{enumerate}

The remainder is organized as follows. Section~\ref{sec:motivation} establishes the empirical motivation (semantic structure, distribution shift). Section~\ref{sec:methodology} describes the integrated system. Sections~\ref{sec:experiments}--\ref{sec:results} present experimental validation with the Pareto frontier as the primary benchmark. Section~\ref{sec:related} surveys related work; Section~\ref{sec:conclusion} discusses limitations and future directions. Extended results and sensitivity analyses appear in the appendices.
